% Automatisch aus der Word-/LaTeX-Konvertierung {\"u}bernommen.
\chapter{Verwandte Arbeiten}

Im dritten Kapitel wird der aktuelle Stand der Forschung zum Thema RTOS-Benchmarking dargestellt und diese Arbeit darin eingeordnet. Zuerst wird erl{\"a}utert, welche Herausforderungen sich beim Vergleich und der Bewertung von Echtzeitsystemen ergeben. Darunter werden auch Fragen zur Reproduzierbarkeit und Vergleichbarkeit zu anderen Betriebssystemen wie FreeRTOS, RIOT oder Contiki-NG bzgl. Speicherverbrauch und Laufzeitverhalten untersucht. Danach werden Zephyr-spezifische Arbeiten behandelt, darunter Zephyrs eigene Referenz-Benchmarks und Studien, die Zephyr in bestimmten Anwendungsszenarien, wie bspw. die Absicherung durch TLS/ DTLS oder KI-Inferenz, einsetzen und teilweise messen. Zuletzt wird diese Arbeit von den vorgestellten Studien abgegrenzt, indem analysiert wird, worin sich ihre Methodik, also der Vergleich von Konfigurationsstufen innerhalb desselben RTOS auf einer spezifizierten Zielplattform, von den anderen bestehenden Vergleichsstudien unterscheidet.

\section{RTOS-Benchmarking allgemein}
Das Benchmarking von Echtzeitbetriebssystemen besteht in der Literatur aus einer Kombination von Laufzeit- und Speichermetriken. Silva et al. (2019) definieren f{\"u}r den Vergleich von Contiki-NG, RIOT und Zephyr vier Bewertungsdimensionen, darunter Leistung, Zeit-Determinismus, Speicher-Fu{\ss}abdruck und Energieverbrauch. F{\"u}r die Leistungs-Bewertung benutzen Silva et al. die Thread-Metric Benchmark Suite. Sie besteht aus acht Tests, darunter kooperative und pr{\"a}emptive Kontextwechsel, Interrupt-Verarbeitung, Semaphor-Operationen, Speicherallokation und -freigabe (Silva et al., 2019, S. 10379). Die Zeitvorhersagbarkeit untersuchen die Autoren zudem anhand eigener Mikrobenchmarks f{\"u}r ausgew{\"a}hlte Thread-Verwaltungsfunktionen. Dar{\"u}ber hinaus beziehen sie einen Thread-Metric-Test ein (Silva et al., 2019, S. 10380). Diese Sammlung an Benchmark-Kategorien {\"a}hnelt den in dieser Arbeit gew{\"a}hlten f{\"u}nf Metriken. Hierbei wurde in dieser Arbeit jedoch nicht die Thread-Metric Benchmarks-Suite verwendet, was genauer in Abschnitt 3.4 begr{\"u}ndet wird.

F{\"u}r die Messung des Zeit-Determinismus benutzen Silva et al. (2019, S. 10379) eigenentwickelte Mikrobenchmarks f{\"u}r vier Thread-Verwaltungs-APIs, darunter \glqq{}Thread Create\grqq{}, \glqq{}Resume\grqq{}, \glqq{}Suspend\grqq{} und \glqq{}Delete\grqq{}. Dabei protokollieren sie Minimum, Maximum und Mittelwert als Jitter-Indikator. Auch Leotescu und Vlad (2018) benutzen Mikrobenchmarks in ihrem Ansatz. Sie entwickeln mit \glqq{}heap\_bench\grqq{}, \glqq{}thread\_bench\grqq{}, \glqq{}mutex\_bench\grqq{} und \glqq{}context\_switch\grqq{} vier eigene, in C geschriebene Testprogramme, die jeweils eine einzelne Kernel-Operation isoliert messen (Leotescu \& Vlad, 2018, S. 3-7). Methodisch {\"a}hnelt dies hierbei dem Aufbau dieser Arbeit, wobei hier ebenfalls jede Metrik in einer eigenen, isolierten \lstinline[style=fileinline]|bm_*.c|-Datei gemessen wird (siehe Kapitel 4.1). Au{\ss}erdem benutzen Leotescu und Vlad (2018, S. 3) f{\"u}r die Zeitmessungen unter Zephyr bewusst die Funktion \lstinline[style=fileinline]|TIMING_INFO_PRE_READ()| anstatt der POSIX-Funktion \lstinline[style=fileinline]|clock_gettime()|, mit der Begr{\"u}ndung, dass die POSIX-Funktion aufgrund des Zephyr-Schedulers inkonsistente Werte liefere. Auch diese Arbeit benutzt mit \lstinline[style=fileinline]|timing_counter_get()| eine Zephyr-spezifische Zeitmessschnittstelle anstatt einer POSIX-Zeitabfrage. Dies entspricht der allgemeinen Entscheidung von Leotescu und Vlad f{\"u}r eine hardwarenahe Zephyr-Zeitbasis. Ihr dabei verwendetes Makro \lstinline[style=fileinline]|TIMING_INFO_PRE_READ()| soll jedoch mit der hier eingesetzten API unterschieden werden (Leotescu \& Vlad, 2018, S. 3).

Jaskani et al. (2019) erg{\"a}nzen dabei diese beiden messtechnisch orientierten Arbeiten um eine breitere, qualitative Einordnung. Sie erw{\"a}hnen n{\"a}mlich sieben Auswahlkriterien f{\"u}r ein IoT-Betriebssystem, darunter Architektur, Programmiermodell, Scheduling-Strategie, Speicher-Fu{\ss}abdruck, Vernetzung, Portabilit{\"a}t und Energieeffizienz (Jaskani et al., 2019, S. 3-4). Zudem ordnen sie Zephyr darin als Mikrokernel-/ Nanokernel-Hybrid mit pr{\"a}emptiven, priorit{\"a}tsbasierten und EDF-f{\"a}higem Scheduler ein (Jaskani et al., 2019, S. 5). Diese Einordnung wird jedoch in dieser Arbeit nicht selbst empirisch nachgepr{\"u}ft, da der Fokus hier nicht auf einem Architektur-, sondern Konfigurationsvergleich liegt.

\section{Bestehende RTOS-Benchmarks}
Silva et al. benutzen f{\"u}r die betreffende Leistungs-Auswertung 1000 Stichproben. Damit liegt der Stichprobenumfang in derselben Gr{\"o}{\ss}enordnung wie bei den zeitbezogenen Benchmarks dieser Arbeit, ohne dass daraus schon eine Gleichheit der Messbedingungen folgt (Silva et al., 2019, S. 10379).

F{\"u}r den untersuchten Zephyr-Aufbau berichten die Autoren au{\ss}erdem einen mittleren RAM-Bedarf von 52,9 KB und einen Flash-Bedarf von 130,5 KB. Die Werte f{\"u}r Contiki-NG betragen 29,8 KB bzw. 50,6 KB und f{\"u}r RIOT 33,0 KB bzw. 59,9 KB. Diese Zahlen basieren dabei auf die jeweils untersuchten Firmware-Aufbauten inklusive Anwendung und Netzwerkfunktionen und nicht auf einen allgemeinen Mindestbedarf der Betriebssysteme (Silva et al., 2019, S. 10381).

Bei den untersuchten Leistungs- und Zeitvorhersagbarkeits-Metriken dokumentieren Silva et al. schlechtere Ergebnisse f{\"u}r Zephyr als f{\"u}r die beiden Vergleichssysteme. Sie diskutieren dabei, dass weitere Verarbeitungs- und Moduswechsel eine m{\"o}gliche Erkl{\"a}rung sein k{\"o}nnen. Auf der anderen Seite notieren sie jedoch, dass der Zeitstempel in ihrem Messaufbau nicht immer am vorgesehenen Endpunkt erfasst wird. Die beobachteten Schwankungen d{\"u}rfen daher nicht unbegrenzt als den realen Jitter der untersuchten API-Operationen interpretiert werden (Silva et al., 2019, S. 10381).

F{\"u}r diese Arbeit ist daraus vor allem die Bedeutung von einer klaren Definition und {\"U}berpr{\"u}fung der Messpunkte relevant. Die Ergebnisse von Silva et al. dienen als Vergleichskontext, ersetzen aber dabei keine Validierung der Zeitmessung, die hier eingesetzt werden.

F{\"u}r diese Arbeit ist diese Erkenntnis auf zwei Arten relevant. Erstens zeigt er, dass Zephyrs Fu{\ss}abdruck und Determinismus-Verhalten stark von der aktivierten Konfiguration abh{\"a}ngen k{\"o}nnen. Genau dieser Zusammenhang wird in dieser Arbeit durch die drei Konfigurationsstufen untersucht. Silva et al. beschreiben die Bedingungen ihres Versuchsaufbaus und benutzen dabei Standard- bzw. empfohlene Einstellungen. Eine komplette Aufstellung aller eingeschalteten Zephyr-Konfigurationssymbole wird jedoch nicht angegeben. Zudem werden innerhalb der Untersuchung unter anderem Thread- und Netzwerkbedingungen variiert. Der Vergleich ist deshalb nicht mit einem gezielten Vergleich von den drei hier definierten Konfigurationsprofile gleichzusetzen (Silva et al., 2019, S. 10378--10380).

Die ver{\"o}ffentlichten Werte machen eine Einordnung der Gr{\"o}{\ss}enordnung m{\"o}glich. Unterschiede in Hardware, Softwarestand, Anwendung und Messverfahren begrenzen dabei jedoch den direkten Vergleich der Zahlen. Daher best{\"a}tigt eine {\"a}hnliche Gr{\"o}{\ss}enordnung allein nicht direkt die Richtigkeit der eigenen Messmethode.

Altoumi und Ahmed beziehen sich dabei UL Solutions Kontextwechselwerte von 223 CPU-Zyklen f{\"u}r FreeRTOS und 524 CPU-Zyklen f{\"u}r Zephyr. Diese Zahlen sind in diesem Kontext keine eigenen Messwerte der Autoren. Sie diskutieren au{\ss}erdem die Bedeutung von unterschiedlichen MPU- bzw. Sicherheitseinstellungen f{\"u}r die Vergleichbarkeit von solchen Ergebnissen (UL Solutions, 2022, zitiert nach Altoumi \& Ahmed, 2025, S. 2--3; Altoumi \& Ahmed, 2025, S. 3).

F{\"u}r diese Arbeit ist daran vor allem der methodische Hinweis relevant, dass Schutzfunktionen bei Zeitmessungen genau dokumentiert werden sollen. Die Minimal-Konfiguration deaktiviert die MPU. Dadurch beziehen sich die Werte, die ermittelt werden, auf eine reduzierte Schutzkonfiguration. Ohne einen weiteren vergleichbaren Messlauf mit aktivierter MPU kann sich der vermiedene Aufwand, der durch diese Abschaltung entsteht, nicht getrennt quantifizieren.

Jaskani et al. (2019) dient in dieser Arbeit an erster Stelle als Kontext- und Einordnungsquelle, nicht als Messreferenz. Das liegt daran, dass das Paper selbst keine eigenen Benchmarks durchf{\"u}hrt. Stattdessen fasst sie bestehende Literatur zu sechs IoT-Betriebssystemen, darunter Contiki, TinyOS, RIOT, Zephyr, Mbed und Brillo zusammen (Jaskani et al., 2019, S. 5-6).

\section{Zephyr-spezifische Benchmarks}
Erg{\"a}nzend zu den in 3.2 vorgestellten RTOS-{\"u}bergreifenden Vergleichsstudien bringt Zephyr selbst eigene Referenz-Benchmarks im Quellbaum mit, der unter \lstinline[style=fileinline]|tests/benchmarks/| liegt. Diese Benchmarks liefern Zephyr-interne Referenzmessungen f{\"u}r ausgew{\"a}hlte Kerneloperationen. In dieser Arbeit werden ihre dokumentierten Messkategorien als Nachvollziehbarkeits- und Abgrenzungspunkte benutzt. Eine unabh{\"a}ngige Validierung der eigenen Messmethodik entsteht daraus allein jedoch noch nicht (Zephyr Project, 2024, tests/benchmarks/latency\_measure/README.rst). Im Folgenden wird f{\"u}r jede der f{\"u}nf untersuchten Metriken notiert, ob ein Zephyr-{\"A}quivalent existiert, was davon {\"u}bernommen und was bewusst anders gemacht worden ist. Erg{\"a}nzend dazu werden drei Studien erw{\"a}hnt, die Zephyr nicht mit anderen RTOS vergleichen, sondern eine Zephyr-Anwendung auf realer Hardware messen.

\textbf{Kontextwechsel-Latenz }F{\"u}r diese Metrik existiert durch \lstinline[style=fileinline]|tests/benchmarks/latency_measure| ein Zephyr-eigenes {\"A}quivalent, das laut der darin enthaltenen \lstinline[style=fileinline]|README|-Datei unter anderem die Kontextwechselzeit {\"u}ber \lstinline[style=fileinline]|k_yield()| misst, sowohl pr{\"a}emptiv als auch kooperativ. Es wird zudem auch die Kontextwechselzeit durch Semaphore-Weckung gemessen. Dieser Benchmark wurde in 5.1 schon als Vergleichspunkt genannt. Wichtig ist hierbei jedoch eine Einschr{\"a}nkung. Der genaue Quellcode von \lstinline[style=fileinline]|latency_measure| wurde im Rahmen dieser Arbeit nicht Zeile f{\"u}r Zeile aus dem GitHub-Repository {\"u}bernommen, sondern nur {\"u}ber die Dokumentation durch die \lstinline[style=fileinline]|README|-Datei eingeordnet. Aus der README werden nur die dort beschriebenen Messkategorien {\"u}bernommen. Die eigene Kalibrierung des Messaufwands ist davon unabh{\"a}ngig entwickelt worden. Die eigene Auswertung f{\"u}hrt Minimum, Maximum, Mittelwert und Stichproben-Standardabweichung laufend {\"u}ber N=1000 Wiederholungen. Sie speichert dabei jedoch keine komplette Verteilung der Einzelwerte.

\textbf{Interrupt-Latenz }Auch dies ist in Zephyrs \lstinline[style=fileinline]|tests/benchmarks/latency_measure| vorhanden. Zephyrs Suite erfasst n{\"a}mlich die Zeit vom Ende der ISR bis zur Fortsetzung eines Threads. Der in dieser Arbeit gemessene Schritt davor, also vom Software-Trigger bis zur ersten ISR, wird von \lstinline[style=fileinline]|latency_measure| jedoch nicht erfasst. Diese Arbeit misst eine Zeitgr{\"o}{\ss}e derselben Gr{\"o}{\ss}enordnung, wobei der Hinweg vom Ausl{\"o}sen bis zur ersten ISR-Instruktion statt des R{\"u}ckwegs gemessen wird. Die README beschreibt f{\"u}r \lstinline[style=fileinline]|latency_measure| den R{\"u}ckweg von einer ISR zum unterbrochenen oder zu einem anderen Thread. Diese Arbeit untersucht auf der anderen Seite den Hinweg vom softwareseitigen Setzen einer NVIC-Anforderung bis zum Zeitstempel am Beginn des eigenen C-Handlers. Aussagen {\"u}ber die bestimmte interne Triggerquelle von \lstinline[style=fileinline]|latency_measure| werden daraus nicht abgeleitet.

\textbf{Timer-Jitter }Im f{\"u}r diese Arbeit gepr{\"u}ften Benchmark-Verzeichnis von Zephyr v3.7 wurde kein Referenzbenchmark gefunden, der die Abweichung der Callback-Intervalle eines periodischen \lstinline[style=fileinline]|k_timer| von seiner Sollperiode misst (Zephyr Project, 2024, tests/benchmarks/).

\textbf{Heap-Allokationsgeschwindigkeit und -Fragmentierung }F{\"u}r Kernel-Operationen gibt es durch \lstinline[style=fileinline]|tests/benchmarks/sys_kernel| ein Zephyr-eigener Benchmark. Dieser misst jedoch nur Semaphore, LIFO, FIFO, Stack und Memory Slab durch \lstinline[style=fileinline]|k_mem_slab_alloc| und \lstinline[style=fileinline]|k_mem_slab_free|. Es wird also nicht der allgemeine System-Heap {\"u}ber \lstinline[style=fileinline]|k_malloc()| und \lstinline[style=fileinline]|k_free()| untersucht. Ein passenderes Zephyr-{\"A}quivalent w{\"a}re daf{\"u}r \lstinline[style=fileinline]|latency_measure| selbst, da es durch \lstinline[style=fileinline]|heap.malloc.immediate| und \lstinline[style=fileinline]|heap.free.immediate| die Heap-Allokation durchf{\"u}hrt. Dabei liefert es, laut der README, in der Standardkonfiguration von Zephyr v3.7.0 Referenzwerte von durchschnittlich 627 Zyklen f{\"u}r \lstinline[style=fileinline]|k_malloc()| und 432 Zyklen f{\"u}r \lstinline[style=fileinline]|k_free()|. F{\"u}r die Heap-Fragmentierung als Zustandsma{\ss} ist kein Zephyr-eigener Benchmark gefunden worden, da \lstinline[style=fileinline]|sys_heap| selbst keine {\"o}ffentliche Kennzahl f{\"u}r den gr{\"o}{\ss}ten zusammenh{\"a}ngenden Freiblock bringt (siehe 5.5). Dies ist auch der Grund f{\"u}r die Eigenentwicklung der Bisektions-Probe in dieser Arbeit.

\textbf{Zephyr in Anwendungsszenarien }Leotescu und Vlad untersuchen vier Mikrobenchmarks f{\"u}r Speicher-, Thread-, Mutex- und Kontextwechsel-Operationen. Dabei vergleichen sie Zephyr auf einem i.MX RT1050 mit Cortex-M7 bei 600 MHz mit Linux auf einem i.MX 6UL mit Cortex-A7 bei 528 MHz. Die verschiedene Hardware begrenzt den direkten Vergleich der Betriebssysteme (Leotescu \& Vlad, 2018, S. 1--3).

F{\"u}r die Zeitmessung unter Zephyr benutzen die Autoren \lstinline[style=fileinline]|TIMING_INFO_PRE_READ()| anstatt \lstinline[style=fileinline]|clock_gettime()|, weil sie mit der POSIX-Zeitabfrage inkonsistente Messwerte beobachtet haben. Dies unterstreicht hierbei die Bedeutung der ausgew{\"a}hlten Zeitbasis. Auch diese Arbeit benutzt f{\"u}r ihre kurzen Laufzeitmessungen eine Zephyr-spezifische Timing-Schnittstelle. Die Eignung muss jedoch f{\"u}r den eigenen Messaufbau beurteilt werden (Leotescu \& Vlad, 2018, S. 3).

In drei wiederholten Benchmark-L{\"a}ufen notieren die Autoren gleiche zusammengefasste Ergebnisse unter Zephyr. F{\"u}r den Kontextwechsel nennen sie einen Mittelwert von 47 Zyklen. Unter Linux beobachten sie Abweichungen von bis zu neun Prozent vom jeweiligen Mittelwert. Diese Ergebnisse gelten dabei f{\"u}r den beschriebenen Versuchsaufbau. Sie zeigen jedoch nicht die Verteilung aller einzelnen Operationen. Au{\ss}erdem zeigen sie auch nicht, dass unabh{\"a}ngige Wiederholungsl{\"a}ufe durch viele Messungen innerhalb eines einzigen Laufs ersetzt werden k{\"o}nnen (Leotescu \& Vlad, 2018, S. 8--9).

Lee et al. untersuchen TLS-/ DTLS-Client-Funktionalit{\"a}t auf Zephyr zuerst in QEMU und anschlie{\ss}end auf einem FRDM-K64F-Board. Das Board benutzt einen Cortex-M4 und verf{\"u}gt {\"u}ber 1 MB Flash und 256 KB RAM. Bei den Hardwaretests notieren die Autoren Kommunikationsprobleme und Paketverluste, die dabei in den vorherigen Simulationstests nicht sichtbar waren (Lee et al., 2018, S. 1293--1294).

Die Arbeit dient hier als Beispiel daf{\"u}r, dass ein erfolgreicher Simulationstest nicht direkt alle Bedingungen von einem realen Ger{\"a}t abbildet. Sie ist dadurch eine Kontextquelle f{\"u}r die Hardware der IoT-Anwendung, aber keine direkte Referenzmessung f{\"u}r die untersuchten Kernel-Latenzen.

Zhang et al. portieren Paddle Lite auf Zephyr und vergleichen die Inferenz auf einem RK3568 mit Cortex-A55 mit einer Linux-Ausf{\"u}hrung. F{\"u}r die untersuchten Modelle notieren sie dabei durchschnittlich sieben Prozent geringere Inferenzzeiten und einen um 78 Prozent geringeren Systemspeicherverbrauch w{\"a}hrend der Inferenz unter Zephyr (Zhang et al., 2026, S. 263--264).

Bei diesen Ergebnissen handelt es sich um einen Anwendungsprozessor und eine KI-Inferenzlast. Sie sind daher nicht direkt auf die Cortex-M4-Plattform und die Kernel-Mikrobenchmarks dieser Arbeit {\"u}bertragbar. Die Studie zeigt jedoch daf{\"u}r, dass Zephyr auch f{\"u}r gr{\"o}{\ss}ere Anwendungen als ressourcenschonende Ausf{\"u}hrungsumgebung untersucht wird.

\section{Abgrenzung der eigenen Arbeit}
Durch die vorgestellten Arbeiten l{\"a}sst sich die Position dieser Arbeit ableiten. Vier Aspekte unterscheiden sie von den genannten Studien.

Erstens vergleicht diese Arbeit nicht, wie die meisten der vorgestellten Studien (Silva et al., 2019; Altoumi \& Ahmed, 2025; Jaskani et al., 2019), unterschiedliche RTOS-Punkte miteinander. Stattdessen wird ein einzelnes RTOS, n{\"a}mlich Zephyr, {\"u}ber drei Konfigurationsstufen hinweg, darunter Minimal, Logging und IoT-Stack, verglichen. Diese Methodik beantwortet daher keinen Produktvergleich wie \glqq{}Welches RTOS ist am schnellsten?\grqq{}, sondern den Einfluss von Konfigurationsentscheidungen auf Laufzeit- und Speicherverhalten desselben Systems. Silva et al. (2019) und Altoumi und Ahmed (2025) bringen zwar interessante Vergleichswerte zwischen RTOS, sagen aber wenig dar{\"u}ber aus, welcher Teil der beobachteten Unterscheide durch die Konfiguration entsteht, anstatt durch das RTOS-Produkt. Diese Arbeit untersucht genau diese L{\"u}cke, indem sie die Konfiguration zur einzigen unabh{\"a}ngigen Variable macht (siehe 4.1).

Zweitens untersucht diese Arbeit Laufzeitmetriken und auch den statischen RAM-/ROM-Fu{\ss}abdruck f{\"u}r die eigenen Zephyr-Konfigurationsstufen auf einer festgelegten Zielplattform. Die Kombination dieser Metrik-Arten ist nicht neu. Silva et al. Achten schon auf Leistung, Zeitvorhersagbarkeit, Speicher und Energieverbrauch. Der weitere Schwerpunkt dieser Untersuchung liegt auf der Variation der Konfiguration innerhalb desselben Betriebssystems (Silva et al., 2019, S. 10378--10382).

Drittens legt diese Arbeit vor allem Wert auf die Nachvollziehbarkeit des eigenen Versuchsaufbaus. Dazu werden die verwendeten Softwarest{\"a}nde, Konfigurationsfragmente, Messpunkte, Kalibrierungsschritte und Auswertungsgr{\"o}{\ss}en dokumentiert. Diese Angaben sollen dabei eine Wiederholung der Untersuchung erleichtern. Die daraus resultierenden Ergebnisse sagen jedoch nicht aus, dass die Dokumentation den Vergleichsstudien {\"u}berlegen ist.

Viertens dienen Beobachtungen aus der Literatur als Ausgangspunkt f{\"u}r eigene Entwurfsentscheidungen. Die Diskussion von unterschiedlichen Schutzkonfigurationen bei Altoumi und Ahmed motiviert bspw., die MPU-Einstellung bewusst zu dokumentieren. Die Ergebnisse zum Speicherbedarf von Silva et al. sind zudem ein Grund f{\"u}r die Untersuchung des Speicherbedarfs von unterschiedlichen Zephyr-Aufbauten. Diese Arbeit {\"u}bernimmt dadurch Fragestellungen aus der Literatur, ohne deren Experimente komplett zu kopieren (Altoumi \& Ahmed, 2025, S. 2--3; Silva et al., 2019, S. 10381).

Insgesamt ist diese Arbeit als eine vertiefte Einzelfallanalyse eines RTOS und kontrollierter Konfigurationsvariation positioniert, die bestehende Literatur-Beobachtungen behandelt und diese auf einer einzigen, genau spezifizierten Hardware-Plattform empirisch nachpr{\"u}ft.
